\section{A corrected Selmer interface and a fixed-coefficient certificate}
\label{sec:sha-odd-bridge}

The earlier report \texttt{grh\_beyond\_dirichlet\_2026-11-25}, in
\texttt{sections/hecke\_selmer.tex}, identifies a finite-precision
quantifier obstruction. Its description matches the displayed theorem
in the compact edition, but it is not an obstruction to the readable
edition's endpoint. The distinction changes the proposed research
program: a full divisible Selmer bound is already present.

\subsection{What the readable edition actually proves}

In \texttt{New/Congruences/Selmer/Line.lean}, lines 120--144,
\texttt{divisibleSelmer} has coefficients
\(\operatorname{Frac}(\mathcal O')/\mathcal O'\), and
\texttt{TwistedLengthLe} bounds the length of this module for all
sufficiently large coefficient rings. The theorem
\texttt{exists\_anticyclotomic\_twist} in
\texttt{GeneralTheorem.lean}, lines 392--396, returns this property for
one character and all its first \(m\) powers. The intermediate precision
\(M\) is chosen with a strict margin above the error bound; it is not a
restriction of the final conclusion to coefficients modulo \(\ell^M\).

The key argument is explicit in
\texttt{Line.lean}, lines 84--95:
\path{length_selmerModule_divisibleQuot_le}. It is useful to expose
its elementary content.

\begin{proposition}[Saturation from one finite level]
\label{prop:one-level-saturation}
Let \(\mathcal O\) be a discrete valuation ring, let \(\pi\) be a
uniformizer, and let \(M\) be a torsion \(\mathcal O\)-module. If
\[
  \operatorname{length}_{\mathcal O}M[\pi^n]<n,
\]
then \(M=M[\pi^n]\), and in particular \(M\) has finite length.
If \(\mathcal O\) is the ring of integers of a finite extension of
\(\mathbb Q_\ell\), with absolute ramification index \(e_{\mathcal O}\),
put
\(\lambda_{\mathcal O}(N)=
\operatorname{length}_{\mathcal O}(N)/e_{\mathcal O}\).
Then
\[
 \lambda_{\mathcal O}\bigl(M[\ell^N]\bigr)\leq b<N
 \quad\Longrightarrow\quad
 M=M[\ell^N],\qquad \lambda_{\mathcal O}(M)\leq b.
\]
\end{proposition}

\begin{proof}
If an element of \(M\) has order \(\pi^r\) with \(r\geq n\), a suitable
multiple generates a submodule isomorphic to \(\mathcal O/\pi^n\)
inside \(M[\pi^n]\). This forces length at least \(n\).
Consequently every element is killed by \(\pi^{n-1}\), so the finite
submodule \(M[\pi^n]\) is all of \(M\).
For the normalized assertion use
\(\ell^N=u\pi^{e_{\mathcal O}N}\) for a unit \(u\).
\end{proof}

The strict inequality is sharp:
\(M=\operatorname{Frac}(\mathcal O)/\mathcal O\) satisfies
\(\lambda_{\mathcal O}(M[\ell^N])=N\) for every \(N\), but has infinite
length.

For the source's Selmer modules, the two algebraic identifications
needed here are proved in the repository. The module of locally
constant cocycles with divisible coefficients is torsion, because
a locally constant function on the compact absolute Galois group
has finite image and hence a common denominator
(\texttt{Line.lean}, lines 62--77). The module with coefficients
\(\mathcal O/a\) identifies with the \(a\)-torsion of the divisible
one (\texttt{Lib/GaloisRep/Main/DivisibleSelmerModules.lean},
lines 118--133). These are modules of cocycles satisfying fixed
linear relations, rather than arbitrary unrelated finite groups.

Thus even the compact theorem has a precise upgrade: retain all its
hypotheses and additionally suppose that its bound satisfies
\(B_i+\eta<M\) for every datum \(i\). After its single twist \(\psi\)
has been chosen, apply Proposition~\ref{prop:one-level-saturation}
to that same twist, coefficient ring, and local relations. The
conclusion is the corresponding full divisible Selmer bound for
\(\psi^d\), \(1\leq d\leq m\). No interchange of
\(\forall M\exists\psi\) and \(\exists\psi\forall M\) occurs.
The readable proof implements this step in
\texttt{twistedLengthLe\_pointChar}, lines 441--470, with
\(v_\ell(P(\zeta-1))<M\); its
\texttt{IsTorsionLevel}, lines 491--499, supplies the required margin.

This corrects one earlier obstruction. It does not identify the
source's local relations with the Bloch--Kato local conditions of an
arbitrary motive, remove its representation hypotheses, or turn a
small normalized length into zero.

\subsection{The exact odd-cocycle functor}

Here is a separate simplification of the remaining noncyclotomic
input. Fix a CM number field \(E\), a prime \(\ell\), a nonnegative
integer \(j\), and a chosen complex conjugation \(c\). Write
\(G_E=\operatorname{Gal}(\overline{\mathbb Q}/E)\), and let
\(\chi_\ell:G_E\longrightarrow\mathbb Z_\ell^\times\) be the cyclotomic
character. For a ring of integers \(\mathcal O\) of a finite extension
of \(\mathbb Q_\ell\), define
\[
 \mathcal A_{\ell,j,\nu}(E;\mathcal O)
\]
to be the \(\mathcal O\)-module of continuous functions
\(z:G_E\to\mathcal O/\ell^\nu\mathcal O\) satisfying:
\begin{enumerate}
\item \(z(gh)=z(g)+\chi_\ell(g)^{-j}z(h)\);
\item \(z\) vanishes on inertia at every finite place away from \(\ell\);
\item if \(j=0\), \(z\) vanishes on inertia above \(\ell\); if \(j>0\),
      \(z\) vanishes on the kernel of \(\chi_\ell\) in every
      decomposition group above \(\ell\);
\item if \(j>0\), \(z(cgc^{-1})=-z(g)\). For \(j=0\) this last
      condition is omitted.
\end{enumerate}
These are exactly the scalar-coordinate conditions of the source's
\texttt{shaOdd}, defined in
\path{Lib/GaloisRep/Main/SelmerModules.lean}, lines 174--190.
No quotient by coboundaries is taken. All local subsets in these
conditions are independent of the coefficient ring and of \(\nu\).

\begin{theorem}[Scalar extension for the source's strict odd cocycles]
\label{thm:sha-basechange}
For every finite coefficient extension \(\mathcal O/\mathbb Z_\ell\),
the natural map is an isomorphism:
\[
 \mathcal O\otimes_{\mathbb Z_\ell}
      \mathcal A_{\ell,j,\nu}(E;\mathbb Z_\ell)
 \ \xrightarrow{\ \sim\ }\
      \mathcal A_{\ell,j,\nu}(E;\mathcal O).
\]
If these modules have finite length, then
\[
 \operatorname{length}_{\mathcal O}
      \mathcal A_{\ell,j,\nu}(E;\mathcal O)
 =
 e_{\mathcal O}\,
 \operatorname{length}_{\mathbb Z_\ell}
      \mathcal A_{\ell,j,\nu}(E;\mathbb Z_\ell).
\]
Consequently their normalized lengths are independent of
\(\mathcal O\) and are nonnegative integers. The statements include
\(\ell=2\).
\end{theorem}

\begin{proof}
Choose a \(\mathbb Z_\ell\)-basis \(b_1,\ldots,b_d\) of \(\mathcal O\).
Reduction modulo the ideal \(\ell^\nu\mathcal O\), not modulo a fixed
power of a uniformizer, gives
\[
 \mathcal O/\ell^\nu\mathcal O
   =\bigoplus_{i=1}^d b_i(\mathbb Z_\ell/\ell^\nu).
\]
Write \(z=\sum_i b_i z_i\). The coordinate maps are maps between finite
discrete coefficient sets, so \(z\) is continuous if and only if all
\(z_i\) are continuous. Each defining equation has coefficients in
\(\mathbb Z_\ell\): this includes the cyclotomic action, both vanishing
conditions, and the equation defining oddness. Hence \(z\) satisfies
the conditions if and only if every \(z_i\) satisfies them. This proves
surjectivity and injectivity of the natural scalar-extension map.
No averaging operator \((1-c)/2\) is used, so no inversion of \(2\)
is required.

For finite length, write the base module as
\(\bigoplus_i\mathbb Z_\ell/\ell^{a_i}\).
After extension it is
\(\bigoplus_i\mathcal O/\ell^{a_i}\mathcal O\), whose length over
\(\mathcal O\) is \(e_{\mathcal O}\sum_i a_i\).
Finite length also descends: if the extended module has finite
length, its restriction to \(\mathbb Z_\ell\) is a finite-length module
containing the base module as a direct summand. The residue degree
does not enter the displayed composition-length factor.
\end{proof}

The earlier \texttt{ramification\_budget.tex} already proved the
generic invariance of normalized length under scalar extension and
identified fixed-module descent as a sufficient extra hypothesis.
Theorem~\ref{thm:sha-basechange} supplies the actual descent
identification for this particular source functor. It does not
supply a comparison from an untwisted Selmer group to a differently
twisted one.

\subsection{A single precision suffices for a proposed odd-cocycle bound}

Let \(\mathcal A_{\ell,j,\infty}(E)\) denote the same cocycle functor
with coefficients in the discrete module
\(\mathbb Q_\ell/\mathbb Z_\ell\).
A continuous function from \(G_E\) to this discrete module has finite
image, because \(G_E\) is compact. Its values therefore have a common
\(\ell\)-power denominator. The natural equivariant embedding
\(\mathbb Z_\ell/\ell^\nu\to\mathbb Q_\ell/\mathbb Z_\ell\),
\(\bar a\mapsto a/\ell^\nu\), gives
\[
 \mathcal A_{\ell,j,\nu}(E;\mathbb Z_\ell)
 \simeq \mathcal A_{\ell,j,\infty}(E)[\ell^\nu],
 \qquad
 \mathcal A_{\ell,j,\infty}(E)
   =\bigcup_{\nu\geq0}\mathcal A_{\ell,j,\infty}(E)[\ell^\nu].
\]
Every defining relation is preserved in both directions. In particular,
the strict local condition has not been replaced by a larger
cohomological condition.

\begin{samepage}
\begin{corollary}[One-precision certificate]
\label{cor:sha-one-precision}
For an integer \(B\geq0\), the following conditions are equivalent:
\begin{enumerate}
\item \(\operatorname{length}_{\mathbb Z_\ell}
        \mathcal A_{\ell,j,\nu}(E;\mathbb Z_\ell)\leq B\)
      for every \(\nu\);
\item \(\operatorname{length}_{\mathbb Z_\ell}
        \mathcal A_{\ell,j,B+1}(E;\mathbb Z_\ell)\leq B\);
\item \(\mathcal A_{\ell,j,\infty}(E)\) has length at most \(B\);
\item for every finite coefficient extension \(\mathcal O\) and every
      \(\nu\),
      \(\lambda_{\mathcal O}
          (\mathcal A_{\ell,j,\nu}(E;\mathcal O))\leq B\).
\end{enumerate}
In particular, the vanishing of the mod-\(\ell\) group is equivalent
to vanishing at every torsion level and over every finite coefficient
extension.
\end{corollary}
\end{samepage}

\begin{proof}
The implication from the first condition to the second is immediate.
The identification with torsion just proved and
Proposition~\ref{prop:one-level-saturation}, with \(N=B+1\), give the
third condition from the second. Taking torsion submodules gives the
first condition from the third.
Theorem~\ref{thm:sha-basechange} proves equivalence with the fourth.
\end{proof}

This is a fixed-field, fixed-exponent, fixed-local-condition
statement. It allows the coefficient ring to vary because a genuine
scalar-extension theorem has been proved. It allows the torsion level
to vary because all truncations belong to the same torsion module.
Changing the character, the representation, the field, or the local
relations requires a separate comparison.

\subsection{The exact remaining arithmetic target}

The source theorem \texttt{exists\_length\_shaOdd\_le}
(\texttt{Lib/Cyclotomic/OddSha.lean}, lines 645--653) requires
\(E\subseteq\mathbb Q(\mu_N)\). It constructs an integer
\[
  0<X\leq 2^{\,4(N+t+3)^{27}}h_E
\]
such that, for every prime \(\ell\), every allowed coefficient ring,
every \(j<t\), and every \(\nu\),
\(\operatorname{length}_{\mathcal O}\mathcal A_{\ell,j,\nu}
  \leq v_{\mathcal O}(X)\), with \(v_{\mathcal O}\) the
integer-valued discrete valuation.

For a noncyclotomic CM field, a proposed positive integer \(X\)
satisfies this entire bound if and only if
\[
 \boxed{\begin{gathered}
 \operatorname{length}_{\mathbb Z_\ell}
    \mathcal A_{\ell,j,v_\ell(X)+1}(E;\mathbb Z_\ell)
       \leq v_\ell(X)\\
 \text{for every prime }\ell\text{ and every }j<t.
 \end{gathered}}
\]
This follows from Corollary~\ref{cor:sha-one-precision}. The
coefficient-ring and all-precision quantifiers have been removed.
The infinitely many prime conditions have not: for every
\(\ell\nmid X\), the condition is mod-\(\ell\) vanishing.
Nor does this result give an algorithm that computes those Galois
cocycle groups.

Finally, the source uses the resulting odd-class cost through
\[
 \sum_\ell v_\ell(X)\log\ell=\log X.
\]
In \texttt{TraceRelations.lean}, lines 106--130, put
\(G=\texttt{chosenGroupOrder}\), \(h=\texttt{halfGap}\), and
\(f=\texttt{numInfPlaces}\). Its denominator error is
\[
 d_\ell=t\,v_\ell(X)
 +\bigl(8ft+2+(4t+1)h^{22}+h\bigr)v_\ell(G),
\]
and its clearing bound is
\[
 \beta_\ell=
   \operatorname{localDefect}_\ell
   +(1+2t)\operatorname{twistError}_\ell+d_\ell+v_\ell(G).
\]
The weighted sum, together with the \(U(1)\) class-number cost,
must fit within
\[
 \frac{\operatorname{errorUnit}}{mk}
   =\frac{fa(\log b)^3}{mk};
\]
see \texttt{FitTwistsExist.lean}, lines 232--260.
A replacement odd-cocycle theorem must therefore control \(\log X\)
uniformly along the varying fields and weights of a transfer.
Establishing a bound separately for each fixed field is insufficient
without this growth estimate.

The results in this section are proof extractions and explicit
interface clarifications. The saturation mechanism is already
formalized upstream; the scalar-extension theorem and its combined
certificate are proved here on paper. We make no claim that these
elementary module-theoretic facts are new in the literature, and we
have not completed a new Lean formalization of this combined interface.

